\chapter{Discussion and Conclusions}

\hspace{\parindent} In this research project, the author presents a standardised computational pipeline that is capable of accurately extracting activity information from automatically segmented cells in performing a comprehensive analysis of two-photon calcium imaging data of any cell model. The proposed pipeline addresses the issues of movement artefacts during imaging of awake mobile mice, the addition of neuropil contamination due to out-of-focus fluorescence from the surrounding neuropil and neighbouring cells, and the heterogeneity in the cellular properties, morphology and temporal activity of the regions of interest. All of these issues impede the accurate detection and segmentation of cells, leading to inaccurate extraction of their neural activity information. 

Among these issues, the problem of movement artefacts poses the greatest challenge in developing an accurate automated analysis pipeline, as the accuracy of the succeeding stages highly depend on the stability of the image sequence, i.e. the positions of the regions of interest do not change over time. If this issue is not properly addressed, the cells will not be properly detected either using morphology-based or activity-based segmentation approaches, or even both. Detecting the cell boundaries using either of these approaches rely on the mean and correlation images of image sequences. Unstable image sequences, i.e. uncorrected data corrupted with either low or high amount of movement artefacts, result to a blurred and noisy mean and correlation images. The evolution of level set functions during the update phase of the level set method is also affected by the instability of the positions of the ROIs of uncorrected data. 

In this study, the author addressed the movement artefacts by applying a subpixel image registration technique that is capable of correcting not only rigid artefacts caused by the natural uniform motion of the brain, but more importantly, it is capable of correcting non-rigid artefacts that arise from the finite speed of raster scanning and the non-rigid deformations due to mechanical strain and head displacement during \textit{in vivo} imaging of head-fixed moving mice. Unlike the HMM-based line-by-line correction approach, \textit{NeuroSEE}'s motion correction algorithm is faster and computationally efficient - making it a suitable method for large high-resolution, high-dimensional imaging datasets. This has been successfully demonstrated when this method was applied to a motion artefact-induced \textit{in vivo} imaging dataset from an awake head-fixed moving mice. The motion-corrected data produced a crisp, more focused and noiseless mean image, indicating stability across the frames in the videos. A significant increase in the correlation-with-mean metric of the motion-corrected data was also shown when compared to that of the raw data. This significant increase validates the significant improvement in the stability of the resulting corrected data. 

In order to validate this further, the proposed hybrid cell detection method in the \textit{NeuroSEE} pipeline was implemented to both raw and motion-corrected data and compared their results. It was shown that the performance of the ROI segmentation module in the motion-corrected data outperforms that of the raw data - proving how the motion correction module improves the detection of cells, as well as the extraction of their neural activity information.

To generalise and make the cell detection more flexible and versatile to a wide variety of imaging datasets, the author integrated both morphology- and activity-based approaches. This is based on the intuition that combining these approaches can improve the general appliabillity of ROI segmentation tools by taking advantage of their individual capabilities, i.e. by detecting temporarily inactive cells using the morphology-based approach and detecting non-circular or non-elliptical-looking cells using the activity-based approach. The integration of these approaches requires the use of both the functional imaging (green) and structural marking (red) channels, a property of the proposed pipeline that is not commonly seen in existing tools. 

Accurate cell detection requires proper tuning of parameters. In this study, the $\alpha$ parameter was tuned optimally to define the appropriate selectivity of initialisation. This parameter is tuned in such a way that it is not too high, i.e. too selective and conservative, resulting to some cells not being detected; but at the same time, it must not be set too low such that it causes overestimation of ROIs, which can make the evolution of the active contours much more complicated and time-consuming. 

After detecting the cells using the hybrid approach, the initial cell boundary estimates, which are represented by active contours, are updated by evolving the level set functions. The criteria in evolving these level set functions are based both on the morphology- and activity-based properties of the two imaging channels. The $\lambda$ parameters was also tuned in order to give a better, smoother final cell boundary estimates. 

In order to evaluate the performance of the proposed \textit{NeuroSEE}'s ROI segmentation module, it was tested on a manually labelled dataset containing manual labels of the expected ROIs which can serve as the ground truth. The predicted cells are matched with the ground truth cells to compute its performance metrics. It was then compared to existing activity-based and morphology-based approaches. It has been shown that \textit{NeuroSEE} has clearly outperformed the other methods in detecting and segmenting cells. It obtained a success rate of 75.4\%, a precision of 72.6\%, and a recall value of 78.4\%, as compared to the other methods whose success rates are all less than 70\%. 

\textit{NeuroSEE} has also successfully decontaminated the extracted calcium traces from neuropil artefacts. By subtracting the weighted subregion neuropil signals from the measured calcium signals, By doing so, the signal-to-noise ratio of the decontaminated signal has improved, leading to a more accurate downstream analysis, such as place field mapping.

Finally, in order to validate whether the extracted calcium signals can be used to analyse the dynamics of the place cells in the hippocampus during mouse exploration and to demonstrate whether these cells in the CA1 region can provide a map-like representation of space, a 2D histogram-based place field mapping was performed. The non-negative derivative of the extracted calcium traces were taken to estimate the onset of spiking of action potentials. These spike onsets were used to construct the rate map of individual cells, which when divided by the occupancy map produces their respective place field maps. With the place maps obtained from the collected neural activity information, it can be shown that \textit{NeuroSEE} is capable of providing an accurate analysis of the dynamics of a population of cells that were detected and segmented automatically. Place field estimation can be used to further quantify information encoding and transmission. 

Thus, the proposed pipeline has shown promising results and has successfully demonstrated that efficient artefact correction and the use of a hybrid approach in detecting and segmenting cells lead to a more accurate and generally applicable extraction of neural activity information. 